\chapter{Технологический раздел}

В данном разделе представляются средства реализации, демонстрация работы сервера и тестирование.

\section{Средства реализации}

В качестве целевой операционной системы для работы данного сервера была выбрана Linux Ubuntu~\cite{ubuntu} по следующим причинам:
\begin{itemize}
	\item возможность реализовать механизм prefork благодаря наличию системного вызова \texttt{fork()}~\cite{fork} в коде ядра;
	\item возможность реализовать механизм select благодаря наличию системного вызова \texttt{select()}~\cite{select} в коде ядра;
	\item возможность реализовать механизм межпроцессного взаимодействия~\cite{ryzanova}.
\end{itemize} 

В качестве языка программирования для реализации данного сервера был выбран C~\cite{c} по следующим причинам:
\begin{itemize}
	\item в стандартной библиотеке языка присутствуют типы и структуры данных, выбранные по результатам проектирования;
	\item код ядра целевой операционной системы написан на данном языке.
\end{itemize}

\section{Демонстрация работы сервера}

Пример работы сервера при запуске представлен в листинге~\ref{lst:start.txt}.

\includelistingpretty{start.txt}{bash}{Пример работы сервера при запуске}

\newpage

Пример отправки HEAD-запроса на сервер представлен в листинге~\ref{lst:head_client.txt}.

\includelistingpretty{head_client.txt}{bash}{Пример отправки HEAD-запроса на сервер}

Пример отправки GET-запроса на сервер представлен в листинге~\ref{lst:get_client.txt}.

\includelistingpretty{get_client.txt}{bash}{Пример отправки GET-запроса на сервер}

Пример работы сервера при обработки HEAD-запроса представлен в листинге~\ref{lst:head_server.txt}.

\includelistingpretty{head_server.txt}{bash}{Пример работы сервера при обработки HEAD-запроса}

\section{Тестирование}

Результаты тестирования работы сервера представлены в таблице~\ref{tab:test_results}.

\begin{table}[H]
\centering
\caption{Результаты тестирования работы сервера}
\label{tab:test_results}
\begin{tabularx}{\textwidth}{|X|X|X|}
	\hline
	\textbf{Тест} & \textbf{Ожидаемый результат} & \textbf{Фактический результат} \\
	\hline
	GET запрос к существующему файлу & HTTP 200 с содержимым файла & Успешно \\
	\hline
	GET запрос к несуществующему файлу & HTTP 404 Not Found & Успешно \\
	\hline
	Запрос с неподдерживаемым методом & HTTP 405 Method Not Allowed & Успешно \\
	\hline
	Path traversal атака & HTTP 403 Forbidden & Успешно \\
	\hline
	GET запрос к директории & HTTP 200 с содержимым директории & Успешно \\
	\hline
	URI превышающий максимальную длину & HTTP 414 URI Too Long & Успешно \\
	\hline
	Запрос файла больше максимального размера & HTTP 413 Payload Too Large & Успешно \\
	\hline
	Множественные одновременные подключения & Корректная обработка всех запросов & Успешно \\
	\hline
	Запрос к файлу с разными MIME-типами & Правильное определение Content-Type & Успешно \\
	\hline
	Запрос при максимальной нагрузке & Стабильная работа без сбоев & Успешно \\
	\hline
	Логирование различных событий & Запись сообщений в указанную цель & Успешно \\
	\hline
\end{tabularx}
\end{table}

Все тесты были успешно пройдены.

\section*{Вывод}

В данном разделе были представлены средства реализации, демонстрация работы сервера и тестирование.